% !TeX root = main.tex
\chapter{Introduction}

\begin{multicols}{2}


\textbf{How the book should be used:} This book is a reference for the game mechanics. It is not a tutorial. Mechanics will be stated only once as clearly and unambiguously as possible to ensure consistency and reading speed. The book is divided into chapters that cover different aspects of the game, where the character creation chapter is necessary to understand the rest of the book.

\textbf{Rules vs Guides:} Rules state exactly how a procedure should be followed. Guides are used to instruct about how the game was designed to be played, but allow freedom of interpretation. A guide will say take these factors into account when making a decision and give them roughly this much importance, but it will not say exactly how much. Rules give players agency by providing actions with clear consequences, while guides will do so by providing expectations and suggestions. The game is designed to be played with a GM, who will arbitrate the rules and interpret the guides. The GM is expected to be fair and consistent, but also to use their judgement to make the game more interesting.

\textbf{Skill based system:} Everything a character does can be resolved as a skill test. The skills are designed to cover any actions in combat, movement, their knowledge and craft related proficiencies, stealth ability and social interactions. All of them are designed to have roughly the same value, so that it is worth investing in any of them, but ultimately the most important ones will be defined by the kind of game being played. 

There are no classes nor levels, so any combination of skills can be obtained by any character, but the game has its mechanisms to specialize characters and give them a unique identity. Investing in some skills will delay the acquisition of others, and the game is designed to make it so that a character cannot be good at everything. 

\section{Design Intent}

\textbf{Verisimilitude.} The game mechanics are designed to roughly map to real world physics. Action points signify time, stamina means breath and muscle fatigue, strength can be expected to be roughly equal to what would be possible in the real world. This sets a common understanding of what is possible and what is not. The point is not to be pedantic about physics, but to make players feel free to use their imagination to solve problems creatively without having to game the rules. If it works in the real world, it should work in the game. If it does not, there should be a good reason for it.

Many of the game mechanics are guidelines on how to design challenges for the character's skills. There is no ambition to set tables and systems to calculate skill test difficulties. What the game offers is a grounded approach from which to estimate the difficulty of a task. This preserves gameplay flow, but deterrs the game from being too whimsical. The GM will be the arbiter of this, and the game will provide the tools to make this task interesting and fair to the players.

\textbf{Player skill vs character skill:} Player skill refers to the player's ability to make decisions and describe plausible actions within the game world. Character skill is a value in the character sheet that is used in a test that corresponds to the action that was described. The game encourages players to describe their actions as if they were the character, and rewards them for doing so. Some actions are not very interesting to describe, in which case the player should just say what they want and a skill check will be made. Other actions will have completely different outcomes depending on how they are done. This is where player skill comes into play. It is what encourages co-creation of the story, not just a GM-run simulation. The difficulty of the test will change depending on how things are done.


\textbf{Play with or without magic.} The magic system is designed in such a way that it does not contradict any laws of physics. Magic is conceived merely as telepathic control over natural phenomena. If the physical ingredients are not present, magic is not possible. This design makes it possible to have a game without magic using the same system. A game without it is complete. 

\textbf{Strategy and risk management.} This game is not about kicking the door and seeing what is inside. It is about making an informed decision before going in. There are many ways to approach a situation and not every strategy is meant to be viable. Sometimes, one strategy will dominate another, even if the losing side has better individual stats. Even low skilled characters are meant to remain a threat to almost anyone, given they are properly equipped. The game philosophy is to give information beforehand, and allow players to prepare for it. 

The game is deadly, as situations with pointy objects tend to be. Death can come from a single blow, even for seasoned warriors. However, the game can be forgiving by allowing players to retreat and try again when they notice that they are outmatched. The point is not to paralyse players with fear, but to reward their ability to assess a situation and make informed decisions.

\textbf{Every rule answers another.} Trust the system to be internally consistent. A rule exists to answer a question that arises from another rule. Some rules may seem to be overpowered, but they carry significant drawbacks that can be exploited. An answer to any mechanic exists within the rules, usually more than one. The game becomes better the more players understand the mechanics.

\textbf{Fast to resolve.} This applies especially to combat. The math is meant to be very expressive of the real world, but also very quick to calculate. It is a compromise solution tuned to get the most depth per time spent. The game is designed to be played with pen and paper, and the math is meant to be done in the head. Online tools are available and do speed up the process, but they are not necessary.

% \section{Types of games}

% \textbf{Historical:} The magic system is optional and can be ignored. This game is good for any time period from antiquity to pike and shot. It is not really designed for modern warfare, since much of the mechanics would become obsolete and unimportant in the presence of modern weaponry.

% \textbf{Horror:} Stress mechanics. Characters lose control. Lighting. Time.

% \textbf{War:} Tactical movement of troops. Technological race. Supply and logistics management. Disposable characters. Helpers. 

% \textbf{Survival and adventure:} Exploration mechanics. Management of inventory, sleep, thirst, and hunger.

% \textbf{Intrigue:} Social. Progress in relationships with NPCs. Social skills. Reputation.

\section{Types of combat}

\textbf{Deadly:} Ends when one team is neutralized.

\textbf{Defend:} Ends when time runs out and the target is protected.

\textbf{Defeat Target:} A specific target must be eliminated. One can flee after defeating the main target.

\textbf{Theft:} The objective is to acquire an object. 

\textbf{Escape:} The objective is to evade pursuers or reach a safe place.

\textbf{Stealth:} The objective is to go unnoticed.

\textbf{Escort/Kidnapping:} The objective is to transport a character.

\textbf{Social} Debate, negotiation, seduction, intimidation. 


\end{multicols}